% Source: book pp. 333--345 (PDF 347--359); no solutions.
% \exnote: Axler's small italic comment under an exercise
\section{Bilinear Forms and Quadratic Forms}

\subsection{Bilinear Forms}

A bilinear form on $V$ is a function from $V\times V$ to $\F$ that is linear in
each slot separately, meaning that if we hold either slot fixed then we have a
linear function in the other slot. Here is the formal definition.

\begin{definition}[Bilinear form]\label{ax:9.1}
A \emph{bilinear form} on $V$ is a function $\beta\colon V\times V\to\F$ such
that
\[ v\mapsto\beta(v,u) \quad\text{and}\quad v\mapsto\beta(u,v) \]
are both linear functionals on $V$ for every $u\in V$.
\end{definition}

\marginnote{Recall that the term \textbf{linear functional}, used in the
definition above, means a linear function that maps into the scalar field $\F$.
Thus the term \textbf{bilinear functional} would be more consistent terminology
than \textbf{bilinear form}, which unfortunately has become standard.}%
For example, if $V$ is a real inner product space, then the function that takes
an ordered pair $(u,v)\in V\times V$ to $\ip{u}{v}$ is a bilinear form on $V$.
If $V$ is a nonzero complex inner product space, then this function is not a
bilinear form because the inner product is not linear in the second slot
(complex scalars come out of the second slot as their complex conjugates).

If $\F=\R$, then a bilinear form differs from an inner product in that an inner
product requires symmetry $\bigl[$meaning that $\beta(v,w)=\beta(w,v)$ for all
$v,w\in V\bigr]$ and positive definiteness $\bigl[$meaning that
$\beta(v,v)>0$ for all $v\in V\setminus\{0\}\bigr]$, but these properties are
not required for a bilinear form.

\begin{example}[Bilinear forms]\label{ax:9.2}\emergencystretch=1.5em\leavevmode
\begin{itemize}
\item The function $\beta\colon\F^3\times\F^3\to\F$ defined by
  \[ \beta\bigl((x_1,x_2,x_3),(y_1,y_2,y_3)\bigr)
     = x_1y_2-5x_2y_3+2x_3y_1 \]
  is a bilinear form on $\F^3$.
\item Suppose $A$ is an $n$-by-$n$ matrix with $A_{j,k}\in\F$ in row $j$,
  column $k$. Define a bilinear form $\beta_A$ on $\F^n$ by
  \[ \beta_A\bigl((x_1,\dots,x_n),(y_1,\dots,y_n)\bigr)
     = \sum_{k=1}^n\sum_{j=1}^n A_{j,k}x_jy_k. \]
  The first bullet point is a special case of this bullet point with $n=3$ and
  \[ A = \begin{pmatrix} 0&1&0\\ 0&0&-5\\ 2&0&0 \end{pmatrix}. \]
\item Suppose $V$ is a real inner product space and $T\in\Lin(V)$. Then the
  function $\beta\colon V\times V\to\R$ defined by
  \[ \beta(u,v) = \ip{u}{Tv} \]
  is a bilinear form on $V$.
\item If $n$ is a positive integer, then the function
  $\beta\colon\Poly_n(\R)\times\Poly_n(\R)\to\R$ defined by
  \[ \beta(p,q) = p(2)\cdot q'(3) \]
  is a bilinear form on $\Poly_n(\R)$.
\item Suppose $\varphi,\tau\in V'$. Then the function
  $\beta\colon V\times V\to\F$ defined by
  \[ \beta(u,v) = \varphi(u)\cdot\tau(v) \]
  is a bilinear form on $V$.
\item More generally, suppose that
  $\varphi_1,\dots,\varphi_n,\tau_1,\dots,\tau_n\in V'$. Then the function
  $\beta\colon V\times V\to\F$ defined by
  \[ \beta(u,v) = \varphi_1(u)\cdot\tau_1(v)+\cdots+\varphi_n(u)\cdot\tau_n(v) \]
  is a bilinear form on $V$.
\end{itemize}
\end{example}

A bilinear form on $V$ is a function from $V\times V$ to $\F$. Because
$V\times V$ is a vector space, this raises the question of whether a bilinear
form can also be a linear map from $V\times V$ to $\F$. Note that none of the
bilinear forms in \ref{ax:9.2} are linear maps except in some special cases in
which the bilinear form is the zero map. Exercise~3 shows that a bilinear form
$\beta$ on $V$ is a linear map on $V\times V$ only if $\beta=0$.

\begin{definition}[$V^{(2)}$]\label{ax:9.3}
The set of bilinear forms on $V$ is denoted by $V^{(2)}$.
\end{definition}

With the usual operations of addition and scalar multiplication of functions,
$V^{(2)}$ is a vector space.

For $T$ an operator on an $n$-dimensional vector space $V$ and a basis
$e_1,\dots,e_n$ of $V$, we used an $n$-by-$n$ matrix to provide information
about $T$. We now do the same thing for bilinear forms on $V$.

\begin{definition}[Matrix of a bilinear form, $\Mat(\beta)$]\label{ax:9.4}\emergencystretch=1.5em
Suppose $\beta$ is a bilinear form on $V$ and $e_1,\dots,e_n$ is a basis of
$V$. The \emph{matrix} of $\beta$ with respect to this basis is the $n$-by-$n$
matrix $\Mat(\beta)$ whose entry $\Mat(\beta)_{j,k}$ in row $j$, column $k$ is
given by
\[ \Mat(\beta)_{j,k} = \beta(e_j,e_k). \]
If the basis $e_1,\dots,e_n$ is not clear from the context, then the notation
$\Mat\bigl(\beta,\allowbreak(e_1,\dots,e_n)\bigr)$ is used.
\end{definition}

Recall that $\F^{n,n}$ denotes the vector space of $n$-by-$n$ matrices with
entries in $\F$ and that $\dim\F^{n,n}=n^2$ (see \ref{ax:3.39} and
\ref{ax:3.40}).

\begin{result}[$\dim V^{(2)}=(\dim V)^2$]\label{ax:9.5}
Suppose $e_1,\dots,e_n$ is a basis of $V$. Then the map
$\beta\mapsto\Mat(\beta)$ is an isomorphism of $V^{(2)}$ onto $\F^{n,n}$.
Furthermore, $\dim V^{(2)}=(\dim V)^2$.
\end{result}

\begin{proof}
The map $\beta\mapsto\Mat(\beta)$ is clearly a linear map of $V^{(2)}$ into
$\F^{n,n}$. For $A\in\F^{n,n}$, define a bilinear form $\beta_A$ on $V$ by
\[ \beta_A(x_1e_1+\cdots+x_ne_n,\,y_1e_1+\cdots+y_ne_n)
   = \sum_{k=1}^n\sum_{j=1}^n A_{j,k}x_jy_k \]
for $x_1,\dots,x_n,y_1,\dots,y_n\in\F$ (if $V=\F^n$ and $e_1,\dots,e_n$ is the
standard basis of $\F^n$, this $\beta_A$ is the same as the bilinear form
$\beta_A$ in the second bullet point of Example~\ref{ax:9.2}).

The linear map $\beta\mapsto\Mat(\beta)$ from $V^{(2)}$ to $\F^{n,n}$ and the
linear map $A\mapsto\beta_A$ from $\F^{n,n}$ to $V^{(2)}$ are inverses of each
other because $\beta_{\Mat(\beta)}=\beta$ for all $\beta\in V^{(2)}$ and
$\Mat(\beta_A)=A$ for all $A\in\F^{n,n}$, as you should verify.

Thus both maps are isomorphisms and the two spaces that they connect have the
same dimension. Hence $\dim V^{(2)}=\dim\F^{n,n}=n^2=(\dim V)^2$.
\end{proof}

Recall that $C^{\mathrm{t}}$ denotes the transpose of a matrix $C$. The matrix
$C^{\mathrm{t}}$ is obtained by interchanging the rows and the columns of $C$.

\begin{result}[Composition of a bilinear form and an operator]\label{ax:9.6}
Suppose $\beta$ is a bilinear form on $V$ and $T\in\Lin(V)$. Define bilinear
forms $\alpha$ and $\rho$ on $V$ by
\[ \alpha(u,v) = \beta(u,Tv) \quad\text{and}\quad \rho(u,v) = \beta(Tu,v). \]
Let $e_1,\dots,e_n$ be a basis of $V$. Then
\[ \Mat(\alpha) = \Mat(\beta)\Mat(T) \quad\text{and}\quad
   \Mat(\rho) = \Mat(T)^{\mathrm{t}}\Mat(\beta). \]
\end{result}

\begin{proof}
If $j,k\in\{1,\dots,n\}$, then
\begin{align*}
\Mat(\alpha)_{j,k} &= \alpha(e_j,e_k)\\
  &= \beta(e_j,Te_k)\\
  &= \beta\Bigl(e_j,\sum_{m=1}^n\Mat(T)_{m,k}e_m\Bigr)\\
  &= \sum_{m=1}^n\beta(e_j,e_m)\Mat(T)_{m,k}\\
  &= \bigl(\Mat(\beta)\Mat(T)\bigr)_{j,k}.
\end{align*}
Thus $\Mat(\alpha)=\Mat(\beta)\Mat(T)$. The proof that
$\Mat(\rho)=\Mat(T)^{\mathrm{t}}\Mat(\beta)$ is similar.
\end{proof}

The result below shows how the matrix of a bilinear form changes if we change
the basis. The formula in the result below should be compared to the
change-of-basis formula for the matrix of an operator (see \ref{ax:3.84}). The
two formulas are similar, except that the transpose $C^{\mathrm{t}}$ appears in
the formula below and the inverse $C^{-1}$ appears in the change-of-basis
formula for the matrix of an operator.

\begin{result}[Change-of-basis formula]\label{ax:9.7}
Suppose $\beta\in V^{(2)}$. Suppose $e_1,\dots,e_n$ and $f_1,\dots,f_n$ are
bases of $V$. Let
\[ A = \Mat\bigl(\beta,(e_1,\dots,e_n)\bigr) \quad\text{and}\quad
   B = \Mat\bigl(\beta,(f_1,\dots,f_n)\bigr) \]
and $C=\Mat\bigl(I,(e_1,\dots,e_n),(f_1,\dots,f_n)\bigr)$. Then
\[ A = C^{\mathrm{t}}BC. \]
\end{result}

\begin{proof}
The linear map lemma (\ref{ax:3.4}) tells us that there exists an operator
$T\in\Lin(V)$ such that $Tf_k=e_k$ for each $k=1,\dots,n$. The definition of
the matrix of an operator with respect to a basis implies that
\[ \Mat\bigl(T,(f_1,\dots,f_n)\bigr) = C. \]

Define bilinear forms $\alpha,\rho$ on $V$ by
\[ \alpha(u,v) = \beta(u,Tv) \quad\text{and}\quad
   \rho(u,v) = \alpha(Tu,v) = \beta(Tu,Tv). \]
Then $\beta(e_j,e_k)=\beta(Tf_j,Tf_k)=\rho(f_j,f_k)$ for all
$j,k\in\{1,\dots,n\}$. Thus
\begin{align*}
A &= \Mat\bigl(\rho,(f_1,\dots,f_n)\bigr)\\
  &= C^{\mathrm{t}}\Mat\bigl(\alpha,(f_1,\dots,f_n)\bigr)\\
  &= C^{\mathrm{t}}BC,
\end{align*}
where the second and third lines each follow from \ref{ax:9.6}.
\end{proof}

\begin{example}[The matrix of a bilinear form on $\Poly_2(\R)$]\label{ax:9.8}
Define a bilinear form $\beta$ on $\Poly_2(\R)$ by
$\beta(p,q)=p(2)\cdot q'(3)$. Let
\[ A = \Mat\bigl(\beta,(1,x-2,(x-3)^2)\bigr) \quad\text{and}\quad
   B = \Mat\bigl(\beta,(1,x,x^2)\bigr) \]
and $C=\Mat\bigl(I,(1,x-2,(x-3)^2),(1,x,x^2)\bigr)$. Then
\[ A = \begin{pmatrix} 0&1&0\\ 0&0&0\\ 0&1&0 \end{pmatrix}
   \quad\text{and}\quad
   B = \begin{pmatrix} 0&1&6\\ 0&2&12\\ 0&4&24 \end{pmatrix}
   \quad\text{and}\quad
   C = \begin{pmatrix} 1&-2&9\\ 0&1&-6\\ 0&0&1 \end{pmatrix}. \]
Now the change-of-basis formula \ref{ax:9.7} asserts that
$A=C^{\mathrm{t}}BC$, which you can verify with matrix multiplication using the
matrices above.
\end{example}

\subsection{Symmetric Bilinear Forms}

\begin{definition}[Symmetric bilinear form, $V_{\mathrm{sym}}^{(2)}$]\label{ax:9.9}
A bilinear form $\rho\in V^{(2)}$ is called \emph{symmetric} if
\[ \rho(u,w) = \rho(w,u) \]
for all $u,w\in V$. The set of symmetric bilinear forms on $V$ is denoted by
$V_{\mathrm{sym}}^{(2)}$.
\end{definition}

\begin{example}[Symmetric bilinear forms]\label{ax:9.10}\leavevmode
\begin{itemize}
\item If $V$ is a real inner product space and $\rho\in V^{(2)}$ is defined by
  \[ \rho(u,w) = \ip{u}{w}, \]
  then $\rho$ is a symmetric bilinear form on $V$.
\item Suppose $V$ is a real inner product space and $T\in\Lin(V)$. Define
  $\rho\in V^{(2)}$ by
  \[ \rho(u,w) = \ip{u}{Tw}. \]
  Then $\rho$ is a symmetric bilinear form on $V$ if and only if $T$ is a
  self-adjoint operator (the previous bullet point is the special case $T=I$).
\item Suppose $\rho\colon\Lin(V)\times\Lin(V)\to\F$ is defined by
  \[ \rho(S,T) = \tr(ST). \]
  Then $\rho$ is a symmetric bilinear form on $\Lin(V)$ because trace is a
  linear functional on $\Lin(V)$ and $\tr(ST)=\tr(TS)$ for all
  $S,T\in\Lin(V)$; see \ref{ax:8.56}.
\end{itemize}
\end{example}

\begin{definition}[Symmetric matrix]\label{ax:9.11}
A square matrix $A$ is called \emph{symmetric} if it equals its transpose.
\end{definition}

An operator on $V$ may have a symmetric matrix with respect to some but not all
bases of $V$. In contrast, the next result shows that a bilinear form on $V$ has
a symmetric matrix with respect to either all bases of $V$ or with respect to no
bases of $V$.

\begin{result}[Symmetric bilinear forms are diagonalizable]\label{ax:9.12}
Suppose $\rho\in V^{(2)}$. Then the following are equivalent.
\begin{parts}
\item $\rho$ is a symmetric bilinear form on $V$.
\item $\Mat\bigl(\rho,(e_1,\dots,e_n)\bigr)$ is a symmetric matrix for every
  basis $e_1,\dots,e_n$ of $V$.
\item $\Mat\bigl(\rho,(e_1,\dots,e_n)\bigr)$ is a symmetric matrix for some
  basis $e_1,\dots,e_n$ of $V$.
\item $\Mat\bigl(\rho,(e_1,\dots,e_n)\bigr)$ is a diagonal matrix for some
  basis $e_1,\dots,e_n$ of $V$.
\end{parts}
\end{result}

\begin{proof}
First suppose (a) holds, so $\rho$ is a symmetric bilinear form. Suppose
$e_1,\dots,e_n$ is a basis of $V$ and $j,k\in\{1,\dots,n\}$. Then
$\rho(e_j,e_k)=\rho(e_k,e_j)$ because $\rho$ is symmetric. Thus
$\Mat\bigl(\rho,(e_1,\dots,e_n)\bigr)$ is a symmetric matrix, showing that (a)
implies (b).

Clearly (b) implies (c).

Now suppose (c) holds and $e_1,\dots,e_n$ is a basis of $V$ such that
$\Mat\bigl(\rho,(e_1,\dots,e_n)\bigr)$ is a symmetric matrix. Suppose
$u,w\in V$. There exist $a_1,\dots,a_n,b_1,\dots,b_n\in\F$ such that
$u=a_1e_1+\cdots+a_ne_n$ and $w=b_1e_1+\cdots+b_ne_n$. Now
\begin{align*}
\rho(u,w) &= \rho\Bigl(\sum_{j=1}^n a_je_j,\sum_{k=1}^n b_ke_k\Bigr)\\
  &= \sum_{j=1}^n\sum_{k=1}^n a_jb_k\rho(e_j,e_k)\\
  &= \sum_{j=1}^n\sum_{k=1}^n a_jb_k\rho(e_k,e_j)\\
  &= \rho\Bigl(\sum_{k=1}^n b_ke_k,\sum_{j=1}^n a_je_j\Bigr)\\
  &= \rho(w,u),
\end{align*}
where the third line holds because $\Mat(\rho)$ is a symmetric matrix. The
equation above shows that $\rho$ is a symmetric bilinear form, proving that (c)
implies (a).

At this point, we have proved that (a), (b), (c) are equivalent. Because every
diagonal matrix is symmetric, (d) implies (c). To complete the proof, we will
show that (a) implies (d) by induction on $n=\dim V$.

If $n=1$, then (a) implies (d) because every 1-by-1 matrix is diagonal. Now
suppose $n>1$ and the implication (a) $\implies$ (d) holds for one less
dimension. Suppose (a) holds, so $\rho$ is a symmetric bilinear form. If
$\rho=0$, then the matrix of $\rho$ with respect to every basis of $V$ is the
zero matrix, which is a diagonal matrix. Hence we can assume that $\rho\ne0$,
which means there exist $u,w\in V$ such that $\rho(u,w)\ne0$. Now
\[ 2\rho(u,w) = \rho(u+w,u+w)-\rho(u,u)-\rho(w,w). \]
Because the left side of the equation above is nonzero, the three terms on the
right cannot all equal $0$. Hence there exists $v\in V$ such that
$\rho(v,v)\ne0$.

Let $U=\{u\in V : \rho(u,v)=0\}$. Thus $U$ is the null space of the linear
functional $u\mapsto\rho(u,v)$ on $V$. This linear functional is not the zero
linear functional because $v\notin U$. Thus $\dim U=n-1$. By our induction
hypothesis, there is a basis $e_1,\dots,e_{n-1}$ of $U$ such that the symmetric
bilinear form $\rho|_{U\times U}$ has a diagonal matrix with respect to this
basis.

Because $v\notin U$, the list $e_1,\dots,e_{n-1},v$ is a basis of $V$. Suppose
$k\in\{1,\dots,n-1\}$. Then $\rho(e_k,v)=0$ by the construction of $U$. Because
$\rho$ is symmetric, we also have $\rho(v,e_k)=0$. Thus the matrix of $\rho$
with respect to $e_1,\dots,e_{n-1},v$ is a diagonal matrix, completing the proof
that (a) implies (d).
\end{proof}

The previous result states that every symmetric bilinear form has a diagonal
matrix with respect to some basis. If our vector space happens to be a real
inner product space, then the next result shows that every symmetric bilinear
form has a diagonal matrix with respect to some \emph{orthonormal} basis. Note
that the inner product here is unrelated to the bilinear form.

\begin{result}[Diagonalization of a symmetric bilinear form by an orthonormal
basis]\label{ax:9.13}
Suppose $V$ is a real inner product space and $\rho$ is a symmetric bilinear
form on $V$. Then $\rho$ has a diagonal matrix with respect to some orthonormal
basis of $V$.
\end{result}

\begin{proof}
Let $f_1,\dots,f_n$ be an orthonormal basis of $V$. Let
$B=\Mat\bigl(\rho,(f_1,\dots,f_n)\bigr)$. Then $B$ is a symmetric matrix (by
\ref{ax:9.12}). Let $T\in\Lin(V)$ be the operator such that
$\Mat\bigl(T,\allowbreak(f_1,\dots,f_n)\bigr)=B$. Thus $T$ is self-adjoint.

The real spectral theorem (\ref{ax:7.29}) states that $T$ has a diagonal matrix
with respect to some orthonormal basis $e_1,\dots,e_n$ of $V$. Let
$C=\Mat\bigl(I,(e_1,\dots,e_n),(f_1,\dots,f_n)\bigr)$. Thus $C^{-1}BC$ is the
matrix of $T$ with respect to the basis $e_1,\dots,e_n$ (by \ref{ax:3.84}).
Hence $C^{-1}BC$ is a diagonal matrix. Now
\[ M\bigl(\rho,(e_1,\dots,e_n)\bigr) = C^{\mathrm{t}}BC = C^{-1}BC, \]
where the first equality holds by \ref{ax:9.7} and the second equality holds
because $C$ is a unitary matrix with real entries (which implies that
$C^{-1}=C^{\mathrm{t}}$; see \ref{ax:7.57}).
\end{proof}

Now we turn our attention to alternating bilinear forms. Alternating
multilinear forms will play a major role in our approach to determinants later
in this chapter.

\begin{definition}[Alternating bilinear form, $V_{\mathrm{alt}}^{(2)}$]\label{ax:9.14}
A bilinear form $\alpha\in V^{(2)}$ is called \emph{alternating} if
\[ \alpha(v,v) = 0 \]
for all $v\in V$. The set of alternating bilinear forms on $V$ is denoted by
$V_{\mathrm{alt}}^{(2)}$.
\end{definition}

\begin{example}[Alternating bilinear forms]\label{ax:9.15}\leavevmode
\begin{itemize}
\item Suppose $n\ge3$ and $\alpha\colon\F^n\times\F^n\to\F$ is defined by
  \[ \alpha\bigl((x_1,\dots,x_n),(y_1,\dots,y_n)\bigr)
     = x_1y_2-x_2y_1+x_1y_3-x_3y_1. \]
  Then $\alpha$ is an alternating bilinear form on $\F^n$.
\item Suppose $\varphi,\tau\in V'$. Then the bilinear form $\alpha$ on $V$
  defined by
  \[ \alpha(u,w) = \varphi(u)\tau(w)-\varphi(w)\tau(u) \]
  is alternating.
\end{itemize}
\end{example}

The next result shows that a bilinear form is alternating if and only if
switching the order of the two inputs multiplies the output by $-1$.

\begin{result}[Characterization of alternating bilinear forms]\label{ax:9.16}
A bilinear form $\alpha$ on $V$ is alternating if and only if
\[ \alpha(u,w) = -\alpha(w,u) \]
for all $u,w\in V$.
\end{result}

\begin{proof}
First suppose that $\alpha$ is alternating. If $u,w\in V$, then
\begin{align*}
0 &= \alpha(u+w,u+w)\\
  &= \alpha(u,u)+\alpha(u,w)+\alpha(w,u)+\alpha(w,w)\\
  &= \alpha(u,w)+\alpha(w,u).
\end{align*}
Thus $\alpha(u,w)=-\alpha(w,u)$, as desired.

To prove the implication in the other direction, suppose
$\alpha(u,w)=-\alpha(w,u)$ for all $u,w\in V$. Then $\alpha(v,v)=-\alpha(v,v)$
for all $v\in V$, which implies that $\alpha(v,v)=0$ for all $v\in V$. Thus
$\alpha$ is alternating.
\end{proof}

Now we show that the vector space of bilinear forms on $V$ is the direct sum of
the symmetric bilinear forms on $V$ and the alternating bilinear forms on $V$.

\begin{result}[$V^{(2)}=V_{\mathrm{sym}}^{(2)}\oplus V_{\mathrm{alt}}^{(2)}$]\label{ax:9.17}
The sets $V_{\mathrm{sym}}^{(2)}$ and $V_{\mathrm{alt}}^{(2)}$ are subspaces of
$V^{(2)}$. Furthermore,
\[ V^{(2)} = V_{\mathrm{sym}}^{(2)}\oplus V_{\mathrm{alt}}^{(2)}. \]
\end{result}

\begin{proof}
The definition of symmetric bilinear form implies that the sum of any two
symmetric bilinear forms on $V$ is a symmetric bilinear form on $V$, and every
scalar multiple of any symmetric bilinear form on $V$ is a symmetric bilinear
form on $V$. Also, the zero bilinear form is in $V_{\mathrm{sym}}^{(2)}$. Thus
$V_{\mathrm{sym}}^{(2)}$ is a subspace of $V^{(2)}$. Similarly, the
verification that $V_{\mathrm{alt}}^{(2)}$ is a subspace of $V^{(2)}$ is
straightforward.

Next, we want to show that
$V^{(2)}=V_{\mathrm{sym}}^{(2)}+V_{\mathrm{alt}}^{(2)}$. To do this, suppose
$\beta\in V^{(2)}$. Define $\rho,\alpha\in V^{(2)}$ by
\[ \rho(u,w) = \frac{\beta(u,w)+\beta(w,u)}{2} \quad\text{and}\quad
   \alpha(u,w) = \frac{\beta(u,w)-\beta(w,u)}{2}. \]
Then $\rho\in V_{\mathrm{sym}}^{(2)}$ and $\alpha\in V_{\mathrm{alt}}^{(2)}$,
and $\beta=\rho+\alpha$. Thus
$V^{(2)}=V_{\mathrm{sym}}^{(2)}+V_{\mathrm{alt}}^{(2)}$.

Finally, to show that the intersection of the two subspaces under
consideration equals $\{0\}$, suppose
$\beta\in V_{\mathrm{sym}}^{(2)}\cap V_{\mathrm{alt}}^{(2)}$. If $u,w\in V$,
then \ref{ax:9.16} implies that
\[ \beta(u,w) = -\beta(w,u) = -\beta(u,w) \]
and hence $\beta(u,w)=0$. Thus $\beta=0$. Hence
$V^{(2)}=V_{\mathrm{sym}}^{(2)}\oplus V_{\mathrm{alt}}^{(2)}$ (by
\ref{ax:1.46}).
\end{proof}

\subsection{Quadratic Forms}

\begin{definition}[Quadratic form associated with a bilinear form, $q_\beta$]\label{ax:9.18}
For $\beta$ a bilinear form on $V$, define a function $q_\beta\colon V\to\F$
by $q_\beta(v)=\beta(v,v)$. A function $q\colon V\to\F$ is called a
\emph{quadratic form} on $V$ if there exists a bilinear form $\beta$ on $V$
such that $q=q_\beta$.
\end{definition}

Note that if $\beta$ is a bilinear form, then $q_\beta=0$ if and only if
$\beta$ is alternating.

\begin{example}[Quadratic form]\label{ax:9.19}
Suppose $\beta$ is the bilinear form on $\R^3$ defined by
\[ \beta\bigl((x_1,x_2,x_3),(y_1,y_2,y_3)\bigr)
   = x_1y_1-4x_1y_2+8x_1y_3-3x_3y_3. \]
Then $q_\beta$ is the quadratic form on $\R^3$ given by the formula
\[ q_\beta(x_1,x_2,x_3) = x_1^2-4x_1x_2+8x_1x_3-3x_3^2. \]
\end{example}

The quadratic form in the example above is typical of quadratic forms on
$\F^n$, as shown in the next result.

\begin{result}[Quadratic forms on $\F^n$]\label{ax:9.20}
Suppose $n$ is a positive integer and $q$ is a function from $\F^n$ to $\F$.
Then $q$ is a quadratic form on $\F^n$ if and only if there exist numbers
$A_{j,k}\in\F$ for $j,k\in\{1,\dots,n\}$ such that
\[ q(x_1,\dots,x_n) = \sum_{k=1}^n\sum_{j=1}^n A_{j,k}x_jx_k \]
for all $(x_1,\dots,x_n)\in\F^n$.
\end{result}

\begin{proof}
First suppose $q$ is a quadratic form on $\F^n$. Thus there exists a bilinear
form $\beta$ on $\F^n$ such that $q=q_\beta$. Let $A$ be the matrix of $\beta$
with respect to the standard basis of $\F^n$. Then for all
$(x_1,\dots,x_n)\in\F^n$, we have the desired equation
\[ q(x_1,\dots,x_n) = \beta\bigl((x_1,\dots,x_n),(x_1,\dots,x_n)\bigr)
   = \sum_{k=1}^n\sum_{j=1}^n A_{j,k}x_jx_k. \]

Conversely, suppose there exist numbers $A_{j,k}\in\F$ for
$j,k\in\{1,\dots,n\}$ such that
\[ q(x_1,\dots,x_n) = \sum_{k=1}^n\sum_{j=1}^n A_{j,k}x_jx_k \]
for all $(x_1,\dots,x_n)\in\F^n$. Define a bilinear form $\beta$ on $\F^n$ by
\[ \beta\bigl((x_1,\dots,x_n),(y_1,\dots,y_n)\bigr)
   = \sum_{k=1}^n\sum_{j=1}^n A_{j,k}x_jy_k. \]
Then $q=q_\beta$, as desired.
\end{proof}

Although quadratic forms are defined in terms of an arbitrary bilinear form,
the equivalence of (a) and (b) in the result below shows that a
\emph{symmetric} bilinear form can always be used.

\begin{result}[Characterizations of quadratic forms]\label{ax:9.21}
Suppose $q\colon V\to\F$ is a function. The following are equivalent.
\begin{parts}
\item $q$ is a quadratic form.
\item There exists a unique symmetric bilinear form $\rho$ on $V$ such that
  $q=q_\rho$.
\item $q(\lambda v)=\lambda^2q(v)$ for all $\lambda\in\F$ and all $v\in V$, and
  the function
  \[ (u,w)\mapsto q(u+w)-q(u)-q(w) \]
  is a symmetric bilinear form on $V$.
\item $q(2v)=4q(v)$ for all $v\in V$, and the function
  \[ (u,w)\mapsto q(u+w)-q(u)-q(w) \]
  is a symmetric bilinear form on $V$.
\end{parts}
\end{result}

\begin{proof}
First suppose (a) holds, so $q$ is a quadratic form. Hence there exists a
bilinear form $\beta$ such that $q=q_\beta$. By \ref{ax:9.17}, there exist a
symmetric bilinear form $\rho$ on $V$ and an alternating bilinear form $\alpha$
on $V$ such that $\beta=\rho+\alpha$. Now
\[ q = q_\beta = q_\rho+q_\alpha = q_\rho. \]
If $\rho'\in V_{\mathrm{sym}}^{(2)}$ also satisfies $q_{\rho'}=q$, then
$q_{\rho'-\rho}=0$; thus
$\rho'-\rho\in V_{\mathrm{sym}}^{(2)}\cap V_{\mathrm{alt}}^{(2)}$, which
implies that $\rho'=\rho$ (by \ref{ax:9.17}). This completes the proof that (a)
implies (b).

Now suppose (b) holds, so there exists a symmetric bilinear form $\rho$ on $V$
such that $q=q_\rho$. If $\lambda\in\F$ and $v\in V$ then
\[ q(\lambda v) = \rho(\lambda v,\lambda v) = \lambda\rho(v,\lambda v)
   = \lambda^2\rho(v,v) = \lambda^2q(v), \]
showing that the first part of (c) holds.

If $u,w\in V$, then
\[ q(u+w)-q(u)-q(w) = \rho(u+w,u+w)-\rho(u,u)-\rho(w,w) = 2\rho(u,w). \]
Thus the function $(u,w)\mapsto q(u+w)-q(u)-q(w)$ equals $2\rho$, which is a
symmetric bilinear form on $V$, completing the proof that (b) implies (c).

Clearly (c) implies (d).

Now suppose (d) holds. Let $\rho$ be the symmetric bilinear form on $V$ defined
by
\[ \rho(u,w) = \frac{q(u+w)-q(u)-q(w)}{2}. \]
If $v\in V$, then
\[ \rho(v,v) = \frac{q(2v)-q(v)-q(v)}{2} = \frac{4q(v)-2q(v)}{2} = q(v). \]
Thus $q=q_\rho$, completing the proof that (d) implies (a).
\end{proof}

\begin{example}[Symmetric bilinear form associated with a quadratic form]\label{ax:9.22}
Suppose $q$ is the quadratic form on $\R^3$ given by the formula
\[ q(x_1,x_2,x_3) = x_1^2-4x_1x_2+8x_1x_3-3x_3^2. \]
A bilinear form $\beta$ on $\R^3$ such that $q=q_\beta$ is given by
Example~\ref{ax:9.19}, but this bilinear form is not symmetric, as promised by
\ref{ax:9.21}(b). However, the bilinear form $\rho$ on $\R^3$ defined by
\begin{multline*}
\rho\bigl((x_1,x_2,x_3),(y_1,y_2,y_3)\bigr)\\
= x_1y_1-2x_1y_2-2x_2y_1+4x_1y_3+4x_3y_1-3x_3y_3
\end{multline*}
is symmetric and satisfies $q=q_\rho$.
\end{example}

The next result states that for each quadratic form we can choose a basis such
that the quadratic form looks like a weighted sum of squares of the
coordinates, meaning that there are no cross terms of the form $x_jx_k$ with
$j\ne k$.

\begin{result}[Diagonalization of quadratic form]\label{ax:9.23}
Suppose $q$ is a quadratic form on $V$.
\begin{parts}
\item There exist a basis $e_1,\dots,e_n$ of $V$ and
  $\lambda_1,\dots,\lambda_n\in\F$ such that
  \[ q(x_1e_1+\cdots+x_ne_n) = \lambda_1x_1^2+\cdots+\lambda_nx_n^2 \]
  for all $x_1,\dots,x_n\in\F$.
\item If $\F=\R$ and $V$ is an inner product space, then the basis in (a) can
  be chosen to be an orthonormal basis of $V$.
\end{parts}
\end{result}

\begin{proof}\leavevmode
\begin{parts}
\item There exists a symmetric bilinear form $\rho$ on $V$ such that
  $q=q_\rho$ (by \ref{ax:9.21}). Now there exists a basis $e_1,\dots,e_n$ of
  $V$ such that $\Mat\bigl(\rho,(e_1,\dots,e_n)\bigr)$ is a diagonal matrix
  (by \ref{ax:9.12}). Let $\lambda_1,\dots,\lambda_n$ denote the entries on the
  diagonal of this matrix. Thus
  \[ \rho(e_j,e_k) = \begin{cases} \lambda_j & \text{if } j=k,\\
     0 & \text{if } j\ne k \end{cases} \]
  for all $j,k\in\{1,\dots,n\}$. If $x_1,\dots,x_n\in\F$, then
  \begin{align*}
  q(x_1e_1+\cdots+x_ne_n)
    &= \rho(x_1e_1+\cdots+x_ne_n,\,x_1e_1+\cdots+x_ne_n)\\
    &= \sum_{k=1}^n\sum_{j=1}^n x_jx_k\rho(e_j,e_k)\\
    &= \lambda_1x_1^2+\cdots+\lambda_nx_n^2,
  \end{align*}
  as desired.
\item Suppose $\F=\R$ and $V$ is an inner product space. Then \ref{ax:9.13}
  tells us that the basis in (a) can be chosen to be an orthonormal basis of
  $V$.\qedhere
\end{parts}
\end{proof}

% ------------------------------------------------------------------ exercises
\begin{exc}

\begin{exercise}
Prove that if $\beta$ is a bilinear form on $\F$, then there exists $c\in\F$
such that
\[ \beta(x,y) = cxy \]
for all $x,y\in\F$.
\end{exercise}

\begin{exercise}
Let $n=\dim V$. Suppose $\beta$ is a bilinear form on $V$. Prove that there
exist $\varphi_1,\dots,\varphi_n,\tau_1,\dots,\tau_n\in V'$ such that
\[ \beta(u,v) = \varphi_1(u)\cdot\tau_1(v)+\cdots+\varphi_n(u)\cdot\tau_n(v) \]
for all $u,v\in V$.
\exnote{This exercise shows that if $n=\dim V$, then every bilinear form on $V$
is of the form given by the last bullet point of Example~\ref{ax:9.2}.}
\end{exercise}

\begin{exercise}
Suppose $\beta\colon V\times V\to\F$ is a bilinear form on $V$ and also is a
linear functional on $V\times V$. Prove that $\beta=0$.
\end{exercise}

\begin{exercise}
Suppose $V$ is a real inner product space and $\beta$ is a bilinear form on
$V$. Show that there exists a unique operator $T\in\Lin(V)$ such that
\[ \beta(u,v) = \ip{u}{Tv} \]
for all $u,v\in V$.
\exnote{This exercise states that if $V$ is a real inner product space, then
every bilinear form on $V$ is of the form given by the third bullet point in
\ref{ax:9.2}.}
\end{exercise}

\begin{exercise}
Suppose $\beta$ is a bilinear form on a real inner product space $V$ and $T$ is
the unique operator on $V$ such that $\beta(u,v)=\ip{u}{Tv}$ for all
$u,v\in V$ (see Exercise~4). Show that $\beta$ is an inner product on $V$ if
and only if $T$ is an invertible positive operator on $V$.
\end{exercise}

\begin{exercise}
Prove or give a counterexample: If $\rho$ is a symmetric bilinear form on $V$,
then
\[ \{v\in V : \rho(v,v)=0\} \]
is a subspace of $V$.
\end{exercise}

\begin{exercise}
Explain why the proof of \ref{ax:9.13} (diagonalization of a symmetric bilinear
form by an orthonormal basis on a real inner product space) fails if the
hypothesis that $\F=\R$ is dropped.
\end{exercise}

\begin{exercise}
Find formulas for $\dim V_{\mathrm{sym}}^{(2)}$ and
$\dim V_{\mathrm{alt}}^{(2)}$ in terms of $\dim V$.
\end{exercise}

\begin{exercise}
Suppose that $n$ is a positive integer and
$V=\{p\in\Poly_n(\R) : p(0)=p(1)\}$. Define $\alpha\colon V\times V\to\R$ by
\[ \alpha(p,q) = \int_0^1 pq'. \]
Show that $\alpha$ is an alternating bilinear form on $V$.
\end{exercise}

\begin{exercise}
Suppose that $n$ is a positive integer and
\[ V = \{p\in\Poly_n(\R) : p(0)=p(1) \text{ and } p'(0)=p'(1)\}. \]
Define $\rho\colon V\times V\to\R$ by
\[ \rho(p,q) = \int_0^1 pq''. \]
Show that $\rho$ is a symmetric bilinear form on $V$.
\end{exercise}
\end{exc}
